\subsection{KT-FF-SEC: finite-field section surjectivity is false}
\label{sec:kt-ff-sec}

Here the prime-to-$p$ arithmetic quotient means
\[
 \Pi_C^{(p')}=\pi_1(C)/
 \ker\bigl(\pi_1(C_{\overline k})\longrightarrow
                   \pi_1(C_{\overline k})^{(p')}\bigr);
\]
the full arithmetic group $G_k$ is retained. This is the convention
of AN\_GRH A05 and of the present question.

\begin{proposition}[KT-FF-SEC]
The map from rational points to conjugacy classes of sections of
$\Pi_C^{(p')}\to G_k$ is not surjective for all proper hyperbolic
curves over finite fields. The same failure occurs for the full
arithmetic fundamental group.
\end{proposition}

\begin{proof}
If $k=\mathbf F_q$, then $G_k\simeq\widehat{\mathbf Z}$.
Every surjection $E\twoheadrightarrow\widehat{\mathbf Z}$ of
profinite groups splits. Indeed, choose $g\in E$ mapping to $1$.
For every open normal subgroup $V\subset E$, the homomorphism
$\mathbf Z\to E/V$ sending $1$ to the image of $g$ factors through
a finite cyclic quotient of $\mathbf Z$. These homomorphisms give
a compatible continuous homomorphism
$s:\widehat{\mathbf Z}\to E$. Its composite with the given
surjection fixes the dense subgroup $\mathbf Z$, and hence is the
identity. This familiar finite-field observation is also the opening
paragraph of Stix, \emph{Rational Points and Arithmetic of Fundamental
Groups}, Chapter 15, p.~197.

For an explicit curve with no rational point, let $C/\mathbf F_2$
be the smooth projective curve with function field
\[
 \mathbf F_2(x,y),\qquad
 y^2+y=1+\frac{x^2+x}{x^3+x+1}.                         \tag{FF.1}
\]
The cubic $x^3+x+1$ has no root in $\mathbf F_2$ and is separable.
Its three roots over $\overline{\mathbf F}_2$ are simple poles of
the right-hand side: none is $0$ or $1$. The rational function has
no other poles. A function of the form $h^2+h$ has even pole order
at every pole, so (FF.1) defines a geometrically connected
Artin--Schreier cover of $\mathbf P^1$ of degree two. A simple pole
has different exponent two; the exact standard reference is
H.~Stichtenoth, \emph{Algebraic Function Fields and Codes}, second
edition, GTM~254 (2009), Proposition~3.7.8(c)--(d), p.~127.
Riemann--Hurwitz therefore gives
\[
 2g(C)-2=2(-2)+3\cdot2=2,
 \qquad g(C)=2.
\]
At each of the three rational points $x=0,1,\infty$ of
$\mathbf P^1(\mathbf F_2)$, the right-hand side is regular and
has value $1$. The corresponding fibre is
$y^2+y=1$, which has no $\mathbf F_2$-solution. Every rational
point of $C$ would map to one of these three rational points.
Thus $C(\mathbf F_2)=\varnothing$.

The fundamental exact sequence of SGA~1, Expos\'e IX,
Th\'eor\`eme 6.1, p.~195 in the 2003 recomposition, gives a
surjection $\pi_1(C)\to G_{\mathbf F_2}$. Its quotient
$\Pi_C^{(2')}\to G_{\mathbf F_2}$ is also surjective. The
preceding splitting argument applies to both, while their
rational-point domain is empty. This proves the claim. Since
$C$ is proper, no missing cuspidal or tangential section can
repair the assertion.
\end{proof}

\paragraph{Scope and attribution.}
The existence of sections over a finite field is standard, not a
new theorem. KT-FF-SEC supplies an explicit application answering
item~6(a) negatively. It does not conflict with the point-injectivity
statement AN\_GRH A05 at cote~149, scans~35--36. Injectivity and
surjectivity are distinct questions. It also does not refute
Grothendieck's section conjecture over finitely generated fields of
characteristic zero, whose base-field hypothesis is absent here.

